\documentclass[11pt]{article}
\usepackage[margin=1in]{geometry}
\usepackage{amsmath,amssymb,amsthm}
\usepackage{booktabs}
\usepackage{algorithm}
\usepackage[noend]{algpseudocode}
\usepackage{listings}
\usepackage{xcolor}
\usepackage[hidelinks]{hyperref}
\usepackage{enumitem}

\lstset{basicstyle=\ttfamily\small,columns=fullflexible,frame=single,
  framerule=0.2pt,breaklines=true,postbreak=\mbox{\(\hookrightarrow\)\space}}

\newtheorem{property}{Property}
\newtheorem{corollary}{Corollary}
\newcommand{\keep}{\mathrm{keep}}
\newcommand{\kept}{\mathrm{kept}}

\title{VestigeKV: An Implementation White Paper\\
\large Training-free sparse attention for NoPE-MLA caches ---\\
a from-zero guide for operator engineers}
\author{VestigeKV project}
\date{2026-09-12 \quad (v2, English; algorithm spec frozen against
\texttt{sglang} port commit \texttt{e44b914a7}, branch \texttt{vestigekv})}

\begin{document}
\maketitle

\begin{abstract}
This document specifies, at implementation depth, the VestigeKV sparse
attention for NoPE-MLA key--value caches: the mathematics of the selection
signal, every data structure, every per-step algorithm, the contract the
fused attention kernel must satisfy, and the validation protocol. It is
written so that an operator engineer holding only this document and stock
model weights can reproduce the method from zero on a new hardware platform
(GPU, NPU, or other accelerators). Everything stated as measured comes from
the pre-registered experimental record of the VestigeKV project; everything
stated as a contract was validated in a production \texttt{sglang} port.
No training, no quantization, no weight or kernel-math changes are involved:
the method is purely a change of \emph{which rows attention reads}.
\end{abstract}

\tableofcontents

%======================================================================
\section{Scope and preconditions}

\subsection{Deliverable scope}
A drop-in replacement for the \emph{decode-time} attention of a NoPE-MLA
transformer layer. Prefill attention, all projections, all other layers, and
the sampling head are untouched. The replacement attends, per decode step,
over a small selected subset of the KV cache instead of all $T$ rows:
\[
\mathrm{attended}(t) \;=\;
\underbrace{\kept(t)}_{\text{rebalanced at each block close}}
\;\cup\;
\underbrace{\mathrm{tail}(t)}_{\text{unclosed rows, incl.\ decoded so far}}
\;\cup\;
\underbrace{\mathrm{recalled}(t)}_{\text{standard, per-step}}
\]
with $|\kept| \approx \rho T$ ($\rho = 1/32$ recommended). Evicted rows are
\emph{never deleted}: they move, bit-exact, to an archive reachable by the
recall tier (Section~\ref{sec:tier2}; the standard configuration, necessary at long context). The deployed invariant is a
\emph{partition} of the cache, not a truncation --- and the partition is
only meaningful if the archive is consulted, so \textbf{the recall tier is
required implementation scope}: a port that implements tier 1 alone is an
incomplete state of this method, not a configuration of it, and must not
expose a switch that disables recall (the only recall knob is the explicit
fetch cap).

\subsection{Architectural preconditions}
\begin{enumerate}[itemsep=1pt]
\item \textbf{MLA (multi-head latent attention) cache.} The per-token cache
row is a single latent vector shared by all heads (MQA-style single KV
head). For Kimi Linear 48B: row width $576 = 512 + 64$, layout
\[
\tilde C_u = [\; \hat c_u \in \mathbb{R}^{512} \;|\; r_u \in \mathbb{R}^{64} \;],
\]
where $\hat c_u$ is the RMS-normalized content latent (\emph{also the value}:
$V$ is read from the first 512 dims) and $r_u$ is the decoupled
``rope branch''.
\item \textbf{NoPE: the branch is never rotated.} In the target model the
rotary call is absent (\texttt{mla\_use\_nope} in the config); $r_u$ is
cached as produced. This is the load-bearing precondition: under RoPE the
identical operator collapses (measured $0.89 \to 0.08$ needle retrieval at
$32\times$), because the selection signal of Section~\ref{sec:signal} exists
only when training has repurposed the positional channel. \emph{A port onto
a RoPE model is unsupported; the collapse is a measured result, not a
hypothesis.} The backend enforces this itself: registration refuses any MLA
config that is not NoPE (and refuses \texttt{page\_size}${\neq}1$ and
speculative decoding), raising rather than degrading silently.
\item \textbf{bf16 cache rows, never quantized.} Both tiers keep the
576-dim latents bit-exact, so every quality effect attributes to selection
and scheduling (a controlled-variable rule). The scan's \emph{index
metadata} is stored 16-bit (side bf16, sketch fp16), calibrated after
rounding; a port that quantizes the cache rows themselves has left the
validated envelope.
\end{enumerate}

\subsection{Mathematical soundness of frozen selection under NoPE}
\label{sec:nope-math}
Write one decode step's attention over cache rows $\{\tilde C_u\}_{u<t}$ with
expanded query $q_t$ (after the standard MLA absorb):
\[
o_t \;=\; \sum_{u} \mathrm{softmax}_u\!\big( s\, q_t \cdot \tilde C_u \big)\, \hat c_u .
\]
Under NoPE both the score map $u \mapsto q_t \cdot \tilde C_u$ and the value
map $u \mapsto \hat c_u$ depend on $u$ \emph{only through the row content}
$\tilde C_u$ --- no term depends on the position index itself.

\begin{property}[Permutation/multiset invariance]
The NoPE-MLA attention output is a function of the \emph{multiset}
$\{\!\{\tilde C_u\}\!\}$ of cache rows (plus the query). Reordering,
re-timing, or re-addressing rows never changes $o_t$.
\end{property}

\begin{corollary}[Ranking stationarity]\label{cor:stat}
Any selection rule that is a function of row contents only can be computed
at \emph{any} time after the rows exist and never goes stale: a kept set
chosen at prefill end is exactly the set the same rule would choose at any
later step over the same rows. Frozen rankings are exact, not approximate.
\end{corollary}

Two engineering consequences follow directly and shape the whole
implementation: (a)~selection runs \emph{off} the decode critical path ---
the implementation recomputes it at prefill end and at every block close
(every $4096$ decoded tokens), where the kept quota $m=\mathrm{round}(\rho\cdot\mathrm{closed})$
grows proportionally and the global top-$m$ is rebalanced over all closed
rows (rows may move in either direction: a long-kept row can be demoted to
the archive); (b)~between closes the per-step sparse pattern is a
\emph{constant} of the decode graph, so it is compatible with
graph-capture execution (CUDA graphs and their equivalents) at zero
per-step host cost. Under RoPE neither holds: scores depend on the
query--key \emph{distance}, so any frozen pattern goes stale at the rate of
the rotary frequency bands.

%======================================================================
\section{The selection signal (tier 1)}\label{sec:signal}

\subsection{Definition}
Selection reads \emph{only} the 64-dim branch slice $r_u$ (11\% of the row;
one contiguous read). The cache prefix is processed in fixed, non-overlapping
blocks of $B = \mathrm{CLOSE\_BLOCK} = 4096$ rows (only full blocks are
scored; the unclosed tail is always attended and needs no score). Within one
block, stack the branches $r \in \mathbb{R}^{B\times 64}$ along the sequence
axis and score each row by its residual against a low-pass reconstruction
along that axis:
\begin{equation}\label{eq:sigma}
\sigma_u \;=\; \big\lVert\, r_u - \big(\mathcal{F}^{-1}\,
\mathbb{1}_{[0,\kappa)}\, \mathcal{F}\, r\big)_u \,\big\rVert_2 ,
\qquad \kappa = 16,
\end{equation}
where $\mathcal{F}$ is the real FFT \emph{within the block} (length $B$,
never the whole prefix), applied independently to each of the 64 channels,
$\mathbb{1}_{[0,\kappa)}$ zeroes all but the lowest $\kappa$ frequency bins,
and $\mathcal{F}^{-1}$ inverts at length $B$. The fixed block window is
load-bearing: $\kappa$ counts frequency bins, and bin $k$ of a length-$B$
transform corresponds to period $B/k$ --- a whole-prefix transform would let
the cutoff drift with context length ($\ge 256$ tokens at $B{=}4096$ vs.\
$\ge 32$k tokens at $T{=}512$k). Per-block transforms keep $\sigma$
scale-invariant, incremental (each block is transformed exactly once, at its
close, and $\sigma$ is immutable thereafter), and bit-identical to the
reference policy. Intuition: under NoPE, training repurposes the positionless
branch into a \emph{salience channel}; rows that carry retrieval-critical
content are spectral anomalies against the smooth background. The signal is
query-independent --- this is what Corollary~\ref{cor:stat} needs.

\subsection{Selection rule}
Over the closed rows only ($\mathrm{closed} = \lfloor T / B \rfloor \cdot B$):
\begin{equation}\label{eq:select}
\keep \;=\; \mathrm{top\text{-}}m(\sigma_{\mathrm{closed}})
\;\cup\; \{0,\dots,\text{sinks}-1\},
\qquad m = \max(1, \mathrm{round}(\rho\cdot\mathrm{closed})),
\end{equation}
with $\text{sinks}=4$ (always keep the first rows). There is \emph{no} fixed
recent window: the unclosed tail $\{\mathrm{closed},\dots,T-1\}$
($<B=4096$ rows, including every token decoded since the last close) is
attended unconditionally and absorbs that role. Top-$m$ is a single
\emph{global} selection over the whole closed prefix, re-run at every block
close with the growing quota $m$ (a proportional-$m$ rebalance, licensed by
Corollary~\ref{cor:stat}). The rule is
applied per layer, independently; only the full-attention (MLA) layers
participate (in Kimi Linear 48B, 7 of 27 layers; the linear-attention layers
have constant-size state and are untouched).

\begin{algorithm}[h]
\caption{Tier-1 selection at prefill end / block close (per MLA layer)}
\begin{algorithmic}[1]
\State $r \gets \text{pool\_rows}[\text{slots of this request}][:, 512{:}576]$
  \Comment{contiguous slice, $[T,64]$}
\For{each full 4096-row block $b$ newly closed since the last selection}
  \State $f \gets \mathrm{rfft}(r[b].\mathrm{float}(), \mathrm{dim}=0)$;
    $f[\kappa{:}] \gets 0$;
    $\ell_b \gets \mathrm{irfft}(f, n{=}4096, \mathrm{dim}=0)$
  \State $\sigma[b] \gets \lVert r[b] - \ell_b \rVert_2 \ \text{over dim} -1$
    \Comment{per-block, immutable once computed}
\EndFor
\State $\keep \gets$ boolean mask from Eq.~\eqref{eq:select} over all closed
  rows \Comment{global top-$m$, $m=\max(1,\mathrm{round}(\rho\cdot\mathrm{closed}))$}
\State write $\mathrm{kept\_buf}[\mathrm{slot}, 0{:}n] \gets$ pool indices of
  $\keep$ rows followed by the unclosed tail (ascending);
  $\mathrm{kept\_len}[\mathrm{slot}] \gets n$
\State invalidate any per-slot decode-side state cached for a previous
  occupant of this slot \Comment{slot-reuse; see Section~\ref{sec:pitfalls}}
\end{algorithmic}
\end{algorithm}

Numerical notes: compute $\sigma$ in fp32 (the branch is bf16 in the pool);
the rFFT length is the block length $4096$, never a padded or whole-prefix
length; the production path is a fused Triton kernel computing the same
statistic in its exact projection form
($\sigma_u = \lVert R_u - C(C^{\top}R)_u \rVert$ over a fixed trigonometric
basis $C$ of the block, fused with a radix-histogram pass for the global
top-$m$) --- mathematically identical to the rFFT chain, ${\sim}28\times$
faster at batched prefill; ties in top-$m$ may break arbitrarily (any
tie-break is valid by Property~1).

\paragraph{Implementation path.} The production path is the fused Triton
kernel \texttt{sigma\_fused}: the rFFT$\to$low-pass$\to$irFFT chain is
evaluated in its exact closed form --- an orthogonal projection onto the
$(2\kappa{-}1)$-dim trigonometric subspace of the fixed block, so
$\sigma_u = \lVert R_u - c_u^{\top}(C^{\top}R)\rVert$ with a precomputed
basis $C$ --- fused with a shared radix histogram (\texttt{atomic\_add} on
the fp32 bit-pattern's high bits) that feeds an exact top-$m$ selection
(\texttt{topm\_from\_hist}: reverse-cumsum the histogram to the boundary bin,
then \texttt{topk} inside it). One launch per block set, ${\sim}28\times$
faster than the batched-\texttt{torch.fft} chain at prefill, and
bit-comparable (fp32-ieee accumulations; the projection form differs from
cuFFT by ${\sim}10^{-7}$ rms against a $1.2{\times}10^{-3}$ bf16 input
quantization floor, and measured top-$m$ agreement is 100\%). The stock
\texttt{torch.fft} / \texttt{torch.topk} chain remains in-tree as the
reference oracle --- a registered test asserts the fused kernel matches it
exactly. Either way this path runs at prefill end and at each block close ---
once per $4096$ generated tokens per request, off the per-token decode path.
Two disciplines are preserved in both forms: (i)~one numerical flavour of
$\sigma$ per block --- never mix two implementations of the same statistic
into one global top-$m$; (ii)~$\sigma$ is computed once per block and never
recomputed. The seven fused Triton kernels of
Section~\ref{sec:count} are the \emph{decode-hot-path} recall step (per
token, per layer); the selection path is a different regime.

%======================================================================
\section{Data structures and memory contract}\label{sec:data}

All structures are per \emph{(request-slot, layer)}; a ``slot'' is the row
index in the runtime's request-to-token table.

\begin{center}
\small
\begin{tabular}{@{}llll@{}}
\toprule
name & shape / dtype & written at & read at\\
\midrule
latent pool & $[N_{\mathrm{pool}}, 576]$ bf16 & every token append & attention kernel\\
$\mathrm{r2t}$ & $[\mathrm{max\_reqs}, \mathrm{ctx}]$ int & runtime-owned & prefill selection\\
$\mathrm{kept\_buf}$ & $[\mathrm{max\_reqs}, \mathrm{ctx}]$ int & prefill; $+1$/step & CSR build\\
$\mathrm{kept\_len}$ & $[\mathrm{max\_reqs}]$ int64 & prefill; $+1$/step & CSR build\\
CSR $(\mathrm{indptr}, \mathrm{indices})$ & $[B{+}1]$, $[\sum n_i]$ int & every decode step & attention kernel\\
\bottomrule
\end{tabular}
\end{center}

\textbf{CSR contract (this is what the kernel consumes).} For batch entry
$i$, the kernel reads rows
$\mathrm{indices}[\mathrm{indptr}[i] \,:\, \mathrm{indptr}[i{+}1]]$ as
absolute pool indices. Boundaries are shared
($\mathrm{end}_i = \mathrm{start}_{i+1}$), so the packed indices must be
contiguous with no gaps. Row order within a request is irrelevant
(Property~1) --- keep ascending order anyway for coalescing.

\textbf{Pad-row convention.} Graph-capture runtimes pad the batch to a
captured bucket size. Every padded CSR slot must point at exactly one
reserved pad row (pool row 0) so its softmax is non-empty; the padded lanes'
outputs are discarded by the runtime. Never let a padded lane read another
request's rows and never leave it an empty range.

\textbf{Growth.} $\mathrm{kept\_len}$ grows by exactly 1 per decode step
(the just-written token's pool slot is appended --- the ``tail''). Size
$\mathrm{kept\_buf}$'s second dim to the context limit; overflow is then
impossible by construction.

\subsection{Index width: int64 today, int32 as a documented lever}
\label{sec:idxwidth}

All slot-id storage ($\mathrm{kept\_buf}$, the fetch buffers, the CSR
indices) is int64 in the shipped implementation. That width is
\emph{inherited}, not needed: the buffers take the dtype of the runtime's
\texttt{kv\_indices} metadata for drop-in compatibility, while the
runtime's own authoritative slot map ($\mathrm{r2t}$) is int32. The value
domain makes int64 pure overhead --- a pool row costs $576\times2$\,B per
MLA layer, so $\mathrm{pool\_rows} < 2\times10^{7}$ on any realistic GPU,
two orders of magnitude below $2^{31}$; int32 cannot overflow.

The cost of the wider word, for Kimi Linear 48B (7 MLA layers) at a
512k context cap with 64 request slots and a batch-32 graph:

\begin{center}
\small
\begin{tabular}{@{}lccc@{}}
\toprule
buffer & shape & int64 & int32\\
\midrule
kept table & $[7,\ 65,\ 524288]$ & 1.9\,GB & 0.95\,GB\\
graph CSR indices & $[7,\ 32\times524288{+}1]$ & 0.94\,GB & 0.47\,GB\\
fetch buffers & $[7,\ 65,\ 4096]$ & 15\,MB & 7\,MB\\
\midrule
total per node & & 2.9\,GB & \textbf{1.4\,GB saved}\\
\bottomrule
\end{tabular}
\end{center}

The compute side of the lever is negligible by design: the recall path
moves only 4-to-8-byte slot ids, $\sim$1/288 of the latent-row traffic, so
halving the index width buys memory, not speed. Conversion is mechanically
simple whenever wanted: \texttt{scatter\_}/\texttt{gather\_} require only
the index \emph{argument} to be int64 (the stored values may be int32, and
the boundary casts already exist), and the Triton pack/scan kernels
promote every loaded index with an explicit \texttt{.to(tl.int64)}, so
they recompile against int32 buffers unchanged. Length and indptr arrays
are kilobyte-scale and stay as they are. We record the lever here rather
than pull it: the reported system keeps the runtime's native index dtype,
and the memory figures in Section~\ref{sec:cost} are therefore
conservative.

%======================================================================
\section{The decode path}\label{sec:decode}

\subsection{Single-request (batch 1): in-graph $O(1)$ append}
With batch 1 the CSR is two scalars, and the per-step state change is one
append. Both are expressible as fixed-shape tensor ops on fixed-address
buffers, so they can be \emph{recorded inside the capture} and replayed with
zero host work:

\begin{algorithm}[h]
\caption{Per-step in-graph refresh, batch 1 (all MLA layers)}
\begin{algorithmic}[1]
\State $\mathrm{slot} \gets \mathrm{req\_pool\_indices}[0{:}1]$;\quad
       $\mathrm{loc} \gets \mathrm{out\_cache\_loc}[0{:}1]$
\For{each local MLA layer $\ell$}
\State $n \gets \mathrm{kept\_len}_\ell.\mathrm{gather}(0, \mathrm{slot})$
\State $\mathrm{flat} \gets \mathrm{slot}\cdot \mathrm{cap} + n$
\State $\mathrm{kept\_buf}_\ell.\mathrm{view}(-1).\mathrm{scatter\_}(0, \mathrm{flat}, \mathrm{loc})$
\State $\mathrm{kept\_len}_\ell.\mathrm{scatter\_}(0, \mathrm{slot}, n{+}1)$
\State $\mathrm{indptr}_\ell[0{:}1] \gets \mathrm{slot}\cdot\mathrm{cap}$;\quad
       $\mathrm{indptr}_\ell[1{:}2] \gets \mathrm{flat}+1$
\EndFor
\end{algorithmic}
\end{algorithm}

The kernel is then bound (at capture time) to
$\mathrm{indices} = \mathrm{kept\_buf}_\ell.\mathrm{view}(-1)$ and
$\mathrm{indptr} = \mathrm{indptr}_\ell$: with batch~1 the request's rows
are already contiguous in its own $\mathrm{kept\_buf}$ row, so no packing is
needed at all.

\textbf{Capture-safety rule (portable).} Inside a capture region, never
extract a scalar from device memory (\texttt{t[0]}-style indexing performs a
hidden device$\to$host sync and invalidates stream capture on CUDA; the
analogous restriction exists on other platforms' graph capture). Use 1-D
gather/scatter/copy with tensor indices only. This restriction is a hard rule
for all graph-capture platforms; the reference port records the corresponding
capture-invalidation failure.

\subsection{Batched decode: out-of-graph vectorized repack}
With batch $>1$, boundaries are shared so a contiguous pack is required each
step. Run it \emph{outside} the captured region (in the runtime's per-step
metadata-refresh hook), fully vectorized --- a per-request loop of small ops
costs milliseconds at batch 16 and was measured to erase the method's gain:

\begin{algorithm}[h]
\caption{Per-step out-of-graph refresh, batch $B$ padded to bucket $B_g$ (per MLA layer)}
\begin{algorithmic}[1]
\State $\mathrm{slots} \gets \mathrm{req\_pool\_indices}[0{:}B]$
  \Comment{$B$ real requests; $\mathrm{out\_cache\_loc}$ has exactly $B$ entries}
\State $n \gets \mathrm{kept\_len}.\mathrm{gather}(0,\mathrm{slots})$;\quad
  append $\mathrm{loc}$ at $(\mathrm{slots}\cdot\mathrm{cap}+n)$;\quad
  $\mathrm{lens} \gets n{+}1$; scatter back
\State $k_{\max} \gets \max(\mathrm{lens})$ \Comment{one small D2H sync}
\State $\mathrm{grid} \gets \mathrm{kept\_buf}.\mathrm{view}(-1).\mathrm{gather}
  \big(0, (\mathrm{slots}[:,\mathrm{None}]\cdot\mathrm{cap} +
  \mathrm{arange}(k_{\max})[\mathrm{None},:])\big)$
\State $\mathrm{valid} \gets \mathrm{arange}(k_{\max})[\mathrm{None},:] < \mathrm{lens}[:,\mathrm{None}]$;\quad
  $\mathrm{total} \gets \mathrm{sum}(\mathrm{lens})$ \Comment{second D2H sync}
\State $\mathrm{indices}[0{:}\mathrm{total}] \gets \mathrm{grid}[\mathrm{valid}]$
  \Comment{masked select = the pack}
\State pad slots $B..B_g{-}1$: one pad row each (row 0);\quad
  $\mathrm{indptr} \gets [0, \mathrm{cumsum}(\mathrm{lens}), \mathrm{total}{+}1, \dots]$
\end{algorithmic}
\end{algorithm}

\textbf{Padding trap (measured crash).} The runtime may pad
$\mathrm{seq\_lens}$/$\mathrm{req\_pool\_indices}$ to the captured bucket
while $\mathrm{out\_cache\_loc}$ keeps only the real entries. Derive the real
batch from $\mathrm{out\_cache\_loc}$'s length, never from the padded shape;
indexing the padded lanes overruns.

\subsection{Prefill and lifecycle}
\begin{enumerate}[itemsep=1pt]
\item \textbf{Short-context dense fallback} (deployment option, default on at
$32\,768$ tokens $=$ 8 close blocks, \texttt{SGLANG\_VESTIGEKV\_ACTIVATION\_MIN\_TOKENS};
$0$ disables): below the threshold selection is off, attended = everything,
byte-identical to stock decode. The threshold exists because short context is
where the method is at a \emph{disadvantage}: the recall pipeline costs a
near-flat $\approx$0.1\,ms/step regardless of context, while dense attention
below $\approx$32k is cheaper than kept+scan+fetch, so compression has
nothing to buy there and only the fixed cost remains. On the crossing step
the whole prefix closes in one pass --- blockwise $\sigma$ is immutable once
computed and the global top-$m$ rebalance commutes with closing order, so the
kept set equals compressing from token 0. \textbf{Every experiment reported
here runs with the fallback disabled} (threshold $0$): the compressed path is
what gets measured at every length, so the short-context numbers below
include the regime where the method pays its fixed cost for no gain --- they
are the honest worst case, not the tuned one.
\item \textbf{Chunked prefill}: run Algorithm~1 at each chunk end (or once at
prefill end); the last write wins and is the frozen kept set.
\item \textbf{Decode}: Algorithms~2/3 per step; the fused kernel runs
unmodified over fewer rows.
\item \textbf{Slot reuse}: a new request occupying a slot must not see any
predecessor state --- invalidate at prefill (see Section~\ref{sec:pitfalls}).
\end{enumerate}

%======================================================================
\section{The attention-kernel contract}\label{sec:kernel}

VestigeKV \emph{reuses} the platform's fused MLA decode kernel; nothing in
the kernel math changes. The port therefore reduces to guaranteeing that the
target platform's kernel satisfies the following contract:

\begin{enumerate}[itemsep=1pt]
\item \textbf{Indexed gather}: keys/values are fetched by absolute pool
indices from the CSR (Section~\ref{sec:data}); rows are \emph{not}
contiguous in the pool. On the validated GPU port the scattered gather cost
is within the measurement noise; on platforms with strong
contiguity requirements, stage the gather through scratch memory.
\item \textbf{MLA absorb decode form}: single shared KV head; per-head
query $q = [q_{\mathrm{nope}}W_{kc} \,|\, q_{\mathrm{rope}}]
\in \mathbb{R}^{576}$; score $= s\, q \cdot \tilde C_u$ with
$s = d^{-1/2}$ per the model config; value $= \tilde C_u[0{:}512]$;
output re-projected by $W_{vc}$ outside the kernel. (This is stock MLA; a platform
already serving DeepSeek-style MLA decode satisfies this item.)
\item \textbf{Split-KV two-stage reduction}: stage 1 computes partial
softmax accumulators + log-sum-exp (LSE) per KV split; stage 2 merges by
LSE. \emph{Required behavior}: empty splits (zero rows) merge out as
$-\infty$ LSE. This is what makes a grid sized for the full context correct
when the CSR shrinks the row count $32\times$: extra splits are empty, not
wrong. Do not resize the grid per step in graph mode; rely on empty-split
correctness.
\item \textbf{Graph capture}: the kernel must read $\mathrm{indptr}$ /
$\mathrm{indices}$ \emph{contents} at execution time from fixed addresses
(the buffers of Section~\ref{sec:data}), so a replayed graph sees each
step's refreshed values. Any metadata precomputed at capture into
registers/immediates breaks the scheme.
\item \textbf{Splice for recall}: the fetched rows join as one extra
split-KV partition (merged by the same LSE path); no kernel change. Graph capture
requires a \emph{fixed fetch-buffer width} $W$, not a cap: the algorithm is
uncapped, and a safety width (default $4096 >$ the observed worst fire of
$3{,}983$) serves its semantics unchanged inside the graph, with overflow
surfaced by telemetry. The cap is merely one policy for choosing $W$.
\end{enumerate}

\textbf{Non-requirements (for tier-1 attention)}: the attention over the
kept+fired rows needs no new kernel, no custom softmax, no changes to
$W_{kc}/W_{vc}$, no rope handling (there is none under NoPE), and no
quantized cache paths. The tier-2 \emph{recall scan} is a separate,
self-contained set of fused kernels (Section~\ref{sec:count}); it touches
the index metadata only, never the attention math or the cache rows.

%======================================================================
\section{The recall tier (tier 2, the standard configuration)}
\label{sec:tier2}
At long context the tier is necessary, not optional: at $S{=}512$k the
reference deployment measures a $99.8\%$ per-step trigger rate with
end-to-end quality passing --- nearly every decode step consults the
archive. A port that stops at tier 1 serves a performance floor, not the
method.

Eviction discards rows whose importance arrives later. The recall tier keeps
every evicted row reachable per step at bounded cost. All thresholds
self-calibrate from the prefix; nothing is trained.

\subsection{The selection principle: certified competition, and why the
fired set is variable}\label{sec:whyvar}

The fired set's size varies per step because selection is not a quota but a
\emph{competition against the query's own attended-tier optimum}. The
complete decision, per head $h$, per archived row $a$:

\begin{enumerate}[itemsep=1pt]
\item \textbf{The standard to beat.} $\mathrm{max}_1^{(h)}$ = the best
attention score this head's query achieves over the \emph{kept} rows. This
is the only quantity an archived row must exceed to matter: a row scoring
below it would receive negligible softmax mass next to what the attended
tier already offers. Everything downstream is query-dependent because this
standard is.
\item \textbf{The gate (skip whole scans).} If the kept-tier softmax is
already confident --- entropy below $\theta_g$ --- the answer is in the
attended tier and the scan is skipped for this head. $\theta_g$ is
self-calibrated as the $\alpha$-quantile of entropy over \emph{hard}
calibration queries (those whose true causal-argmax target is archived),
with $\alpha = (1-\tau)/2 = 5\%$ at $\tau=0.90$, so
at most $\sim$5\% of known-hard queries are wrongly gated off; if that
threshold would pass more than 60\% of all queries the gate cannot separate
and disables itself ($\theta_g = -\infty$).
\item \textbf{The proxy score.} The true score $s\,q\cdot\tilde C_a$ cannot
be afforded against every archived row; the index stores an exact part and
a sketch: $\mathrm{idx}_a = s\,(q_{\mathrm{rope}}\!\cdot\!\mathrm{side}_a +
q_{\mathrm{sk}}\!\cdot\!\mathrm{csk}_a)$ --- exact on the 64-dim branch,
rank-$r$ on the content.
\item \textbf{The certificate (why a proxy can be trusted).} The truncated
part of the true score is $\langle q_\perp, c_\perp\rangle$, the inner
product of the query's and the row's components outside the sketch
subspace. By Cauchy--Schwarz it is at most
$\lVert q_\perp\rVert\,\lVert c_\perp\rVert = q_{\mathrm{res}}\varrho_a$;
scaled by $1/\sqrt{512-r}$ it becomes an expected-magnitude (isotropic)
correction rather than the worst case, and $z$ dials how much of it to
credit: $\mathrm{score}_a = \mathrm{idx}_a + z\,s\,q_{\mathrm{res}}
\varrho_a/\sqrt{512-r}$.
\item \textbf{$z$ self-calibration (the recall guarantee, in closed form).}
The tier has one quality parameter, $\tau$ (default $0.90$); everything else
is derived from it. Split the miss budget $1-\tau$ evenly between the two
stages that can lose a target: the entropy gate may falsely close on a hard
query with probability at most $\alpha = (1-\tau)/2$, so the scan itself must
recover at level $\tau_s = \tau/(1-\alpha)$, giving overall recall
$(1-\alpha)\,\tau_s = \tau$.

For a hard calibration query $q$ the archived target wins exactly when
$\mathrm{idx}_t + z\,\mathrm{cert}_t > \mathrm{max}_1$, i.e.\ when
$z > z_q := (\mathrm{max}_1 - \mathrm{idx}_t)/\mathrm{cert}_t$. Over $n$
exchangeable hard samples, take the order statistics
$z_{(1)} \le \cdots \le z_{(n)}$ and set
\[ z \;=\; z_{(k)}, \qquad k \;=\; \lceil (n{+}1)\,\tau_s \rceil . \]
This is the conformal quantile: by exchangeability alone --- no distributional
assumption on queries or rows --- a fresh hard query's required inflation
falls at or below $z_{(k)}$ with probability at least $\tau_s$, so
$P(\text{target recovered}) \ge \tau_s$ marginally. The quantile exists iff
$n \ge \tau_s/(1-\tau_s)$ (18 samples at $\tau=0.9$), which \emph{derives}
the calibration sample size: collect decode queries, doubling the window,
until that many hard samples exist (or a bounded cap is hit, in which case
$z$ clamps to a safety maximum $z_{\max}$ and errs toward over-fetching,
never under-recall). An earlier implementation searched an 11-rung ladder
for the smallest passing $z$; the ladder was a discretization of this same
quantile that could also undershoot its own target between rungs, and is
gone --- along with the ladder itself, its fallback rung, the fixed
calibration count, and the separate minimum-hard-sample knob, all of which
this one formula replaces.
\item \textbf{Fire and fetch.} Row $a$ fires iff
$\mathrm{score}_a > \mathrm{max}_1^{(h)}$ and the gate is open; the fetch
is the union over heads (per-head top-$j$ capped first when the cap is
set --- capped selection is thus ``the top-$j$ among certified winners'',
keeping the competition's precision while bounding volume).
\end{enumerate}

The variability now has a one-line explanation: \emph{the fired count is
the number of rows that certifiably might beat what the attended tier
already gives this query}. A well-served query fires zero (high
$\mathrm{max}_1$, closed gate); a query whose target was evicted fires the
few rows that plausibly hold it; at $S{=}512$k nearly every step fires
something (the measured $99.8\%$ trigger rate) because at that length the
kept tier alone rarely dominates. A fixed quota cannot reproduce this: set
below a query's need it misses targets, set above it fetches rows that
lose the competition anyway. The variable set is the smallest set
maintaining the calibrated recall target per query; the explicit cap then
trades a bounded slice of that guarantee for a hard latency bound
(Section~\ref{sec:params}).

\subsection{Index build (per layer; provisional at the first decode step,
calibrated asynchronously)}\label{sec:build}
The index is deliberately \emph{not} built at prefill end (under chunked
prefill the last chunk is not reliably identifiable from the backend, and the
calibration queries do not exist yet). The build has two stages:
\begin{enumerate}[itemsep=1pt]
\item \textbf{Provisional build (first decode step, synchronous).} A
conservative index serving from step one: $V$ = the first $r$ rows of the
identity (no SVD yet), gate fully open ($\theta_g=-\infty$), $z = Z_{\max}$
--- over-fetch, never under-recall. (Cache-row proxy queries were tried and
rejected: a row used as its own query has itself as causal argmax, always
inside the attended tail, so \emph{no} sample is ever hard and $z$ flaps
uncalibrated.)
\item \textbf{Calibrated build (asynchronous, off the decode stream).} Real
decode queries are collected per layer (window starts at $8$ steps,
grows by the measured hard rate up to $64$); once enough hard samples exist,
\texttt{RecallTier.build} runs on a side CUDA stream and is installed
atomically. A layer calibrated once donates its basis ($v\_{\text{init}}$)
to later rebuilds --- the certificate is a Cauchy--Schwarz bound, sound for
\emph{any} row-orthonormal basis, so a re-used basis only loosens, never
misses.
\end{enumerate}
Inputs: pool rows $C \in \mathbb{R}^{T\times 576}$ (fp32 for the build), the
tier-1 mask $\keep$, and $n$ \emph{calibration queries}
$q^{\mathrm{cal}} \in \mathbb{R}^{n\times H\times 576}$ (real decode queries
in absorbed form; $H$ heads), with their positions.

\begin{enumerate}[itemsep=1pt]
\item Archive $\mathcal{A} = \{u : \neg\keep_u\}$; $A = |\mathcal{A}|$.
\item \textbf{Sketch basis}: center the calibration queries' content part
and take the top-$r$ eigenvectors of its covariance
($r=64$ recommended GPU-resident):
$V \in \mathbb{R}^{r\times 512} \gets$ rows of
$\mathrm{eigh}\big((q^{\mathrm{cal}}_{:, :512}
- \overline{q^{\mathrm{cal}}_{:, :512}})^{\!\top}
(q^{\mathrm{cal}}_{:, :512} - \overline{q^{\mathrm{cal}}_{:, :512}})\big)$
with the $r$ largest eigenvalues, in descending order. (Equivalently the
top-$r$ right singular vectors; the eigh form is what the code ships ---
measured $17.4 \to 4.3$\,ms at unit subspace alignment, where randomized
\texttt{svd\_lowrank} aligned only $0.87$ and was rejected.)
\item \textbf{Per-row index entries}, for $a \in \mathcal{A}$:
sketch $\mathrm{csk}_a = \hat c_a V^{\!\top} \in \mathbb{R}^{r}$;
residual norm $\varrho_a = \lVert \hat c_a - \mathrm{csk}_a V\rVert$;
branch $\mathrm{side}_a = r_a \in \mathbb{R}^{64}$.
Index footprint per row: $r + 1 + 64$ floats.
\item \textbf{Calibration labels}: for each calibration query, find its
causal argmax row over the \emph{full} cache (chunk the key axis; the pool
is complete by the partition invariant). A query is \emph{hard} iff its
target is archived.
\item \textbf{Calibration queries must be REAL queries.} The gate and $z$
are calibrated on \emph{hard} queries --- those whose causal-argmax target is
archived. Standing in cache rows for queries (an attractive shortcut, since a
row and an expanded query share the 576-dim space in the absorbed form) is
degenerate, not merely biased: a cache row used as its own query has itself as
causal argmax, and its own position always lies inside the recent window, which
tier~1 always keeps. So \emph{no} sample is ever hard, by construction. Measured
on the \texttt{sglang} port: $n_{\mathrm{hard}}=0$ on every build, $z$ falling
back to an uncalibrated constant, and firing on $1.31\%$ of the archive against
the reference record's $0.024\%$ --- a $60{\times}$ inflation of the fetch
volume, invisible in any output, visible only as latency.

Where an engine can reach real prefill queries, use the last
$n_{\mathrm{cal}}$ of them. Where it cannot --- an attention backend sees MLA
prefill on the un-absorbed path, $[\,\cdot\,,H,192]$, and absorbing to the
576-dim index needs $W_{kc}$, a model weight --- calibrate instead on the
request's first $n_{\mathrm{cal}}$ \emph{decode} queries, which arrive already
absorbed. Recall stays inactive for those steps, which is the price. On the
port this restored $n_{\mathrm{hard}}$ to $3$--$63$ per layer, per-layer
calibrated $z$ landing continuously in $[0, Z_{\max}]$, and firing on $0.021\%$ of the archive
--- the reference's regime. Decode queries are arguably the better calibration
set regardless: they are the distribution recall actually serves.
\item \textbf{Entropy gate}: over kept rows only, per query compute
$\mathrm{ent} = H(\mathrm{softmax}(s\,q\,C_{\kept}^{\!\top}))$. Set the gate
threshold $\theta_g$ to the $\alpha$-quantile of $\mathrm{ent}$ over hard
queries, $\alpha=(1-\tau)/2$ (only when at least
$\mathrm{min\_hard}=\lceil\tau_s/(1-\tau_s)\rceil=18$ hard queries exist;
otherwise the gate stays open); if the resulting gate would pass $>60\%$ of
all queries, disable the gate ($\theta_g=-\infty$). The gate exists to skip
scans when the kept-set softmax is already confident.
\item \textbf{Certificate weight $z$ (closed form).}: for hard queries, with
$q_{\mathrm{sk}} = q_{:512}V^{\!\top}$,
$q_{\mathrm{res}} = \lVert q_{:512} - q_{\mathrm{sk}}V\rVert$
(equivalently $\sqrt{\lVert q_{:512}\rVert^2 - \lVert q_{\mathrm{sk}}\rVert^2}$
since $V$ is row-orthonormal), target-row
proxy score
$\mathrm{idx}_t = s\,(q_{512:}\!\cdot\!\mathrm{side}_t +
q_{\mathrm{sk}}\!\cdot\!\mathrm{csk}_t)$ --- computed with the query operands
\emph{rounded to their storage dtypes first} (bf16 side, fp16 sketch), so
calibration matches the serving kernel bit-for-bit (quantize-then-calibrate)
--- and certificate
$\mathrm{cert}_t = s\, q_{\mathrm{res}}\varrho_t / \sqrt{512-r}$,
each hard query requires $z > z_q := (\mathrm{max}_1 - \mathrm{idx}_t) /
\mathrm{cert}_t$; set $z$ to the conformal order statistic
$z = \min\big(z_{(k)},\, Z_{\max}\big)$ with
$k = \lceil (n{+}1)\tau_s \rceil$ over the $n$ hard samples and
$Z_{\max}=8$. If $n < \mathrm{min\_hard}$ (or $n=0$), clamp $z = Z_{\max}$
and keep collecting (over-fetch, never under-recall). ($z$ inflates the
proxy by a Cauchy--Schwarz-style bound on the sketch truncation error, so a
true target is not missed for lack of sketch rank.)
\end{enumerate}

\subsection{Per-step query (each decode step; per layer)}
Given the step's expanded queries $q \in \mathbb{R}^{H\times 576}$:
\begin{align}
\mathrm{max}_1, \mathrm{ent} &\gets \text{from } s\,q\,C_{\kept}^{\!\top}
&&\text{(the kernel's own LSE pass can supply this)}\notag\\
\mathrm{gate} &\gets \mathrm{ent} > \theta_g &&[H]\notag\\
\mathrm{score} &\gets s\big(q_{512:}\mathrm{side}^{\!\top} +
q_{\mathrm{sk}}\,\mathrm{csk}^{\!\top}\big) +
z\, s\,\frac{q_{\mathrm{res}}\,\varrho^{\!\top}}{\sqrt{512-r}}
&&[H,A] \;\;\text{(one GEMV-shaped scan)}\notag\\
\mathrm{fire} &\gets (\mathrm{score} > \mathrm{max}_1) \wedge \mathrm{gate}
&&[H,A]\notag\\
\mathrm{fetch} &\gets \mathrm{fire} \wedge \mathrm{topj\text{-}cap}(\mathrm{score})
&&\text{per head, if } \mathrm{topj}>0\notag
\end{align}
Fetch the union over heads of fired archive rows and splice them into this
step's attention as one extra split-KV partition. Complexity: one
$[H \times A] \cdot (r{+}65)$ scan per layer per step.

\subsection{Fused in-graph implementation (the shipped scan)}
The per-step query above is one GEMV-shaped scan \emph{mathematically};
the shipped kernel realizes the whole recall step --- prologue, scan,
compaction, and CSR pack --- as \textbf{seven fused Triton kernels baked
into the decode CUDA graph itself}, so it costs one \texttt{cudaGraphLaunch}
per step with no second graph and no VestigeKV-triggered recapture:
\begin{enumerate}[itemsep=1pt]
\item \textbf{prologue (2 kernels)}: an online-softmax pass over the kept
rows (forked from the platform's decode-attention accumulation) supplies
$\mathrm{max}_1$ and the entropy gate; a split-NK layout restores occupancy
(one CTA per (pair, head-block)), and the merge kernel folds in the
$q_{\mathrm{sk}}/q_{\mathrm{side}}/q_{\mathrm{res}}$ projections. Queries
are read in place from the stacked \texttt{qbuf}; no materialized gather.
\item \textbf{scan (1 kernel)}: native bf16/fp16 tensor-core dots with
fp32-ieee accumulation over a worst-case-capacity pack; each program covers
one 1024-row compact bucket and emits its fired count as a plain store.
\item \textbf{compaction (2 kernels)}: an exclusive prefix over buckets
(self-zeroing for the next replayed step) then an ordered fired-row write
--- deterministic by construction, so the fired layout is identical across
replays (a bit-parity requirement).
\item \textbf{CSR pack (2 kernels)}: forked from the platform's
kv-indices builder; a single-program-per-layer prep applies the step's
append and writes indptr from pre-step state only (no cross-program
visibility race), then a (layers $\times$ lanes) gather streams kept+fired
rows into the CSR the attention kernel reads.
\end{enumerate}
Correctness rests on three invariants: (i) the pack and every baked buffer
are sized for the worst case at startup, so contents refresh in place and
the model graph never recaptures for VestigeKV reasons; (ii) placeholder
pairs and padded lanes carry zero lengths and point at a trash slot, so the
always-running kernels fire nothing and write where nothing reads --- and
the grid is clipped per batch-size class to the live pair count, so
placeholders are not even launched; (iii) the scan uses the
\emph{previous} step's query (stale-by-one), so it can run at the top of
the captured step before any layer overwrites \texttt{qbuf}. Every kernel
retains its pre-fusion implementation in-tree as a bit-parity oracle, and
the fused pipeline's 2048-token greedy output is token-identical to the
unfused one. A kill-switch env var falls back to a separately captured
scan graph.

\subsection{How the per-step recall count arises}\label{sec:count}
The most-asked implementation question is how the number of recalled rows
is computed. The precise answer: it is not computed --- it \emph{emerges}
as the cardinality of a union of per-head threshold sets,
\[
N_{\mathrm{step}} \;=\; \Big|\, \bigcup_{h\,:\,\mathrm{gate}_h}
\big(\{a \in \mathcal{A} : \mathrm{score}_{h,a} > \mathrm{max}_1^{(h)}\}
\;\cap\; \underbrace{\mathrm{top}\text{-}j_h}_{\text{optional clip}}\big) \Big| ,
\]
produced by a six-stage pipeline per layer per step (shapes in brackets):
(1)~kept-tier scores give per-head $\mathrm{max}_1$ and entropy $[H]$ (the
kernel's LSE pass supplies both); (2)~the gate selects the scanning heads,
$H' = |\{h : \mathrm{ent}_h > \theta_g\}|$ --- the first source of
variation; (3)~one GEMV-shaped scan scores $[H' \times A]$;
(4)~thresholding against each head's own $\mathrm{max}_1$ yields per-head
crossing counts $F_h$ --- the second source; (5)~the cap clips each head to
$\min(F_h, j)$; (6)~the union over heads deduplicates --- the third source
(head overlap).

There is deliberately no closed-form count: $N_{\mathrm{step}}$ is the
measure of the score distribution's level set above a query-dependent
threshold --- data, not configuration. What \emph{is} configured is the
distribution's shape, through exactly two knobs. The calibration target
(via $z$, the conformal quantile of required inflations at the 90\%
recall criterion) sets the operating point, making $N_{\mathrm{step}}$ the
size of the smallest certified set at that point --- one row more is
waste, one fewer breaks the guarantee; raising the target shifts the whole
distribution up. The cap then clips the tail at $j \times H'$.

Measured shape at $S{=}512$k (reference record): trigger rate $99.8\%$ of
steps ($N \ge 1$), mean fetch $0.024\%$ of the archive ($\approx$120 rows),
uncapped worst observation $3{,}983$, capped worst $296$. The distribution
is frequent-and-small with a heavy tail --- which is precisely what the cap
exists to clip, and how a serving SLO should quote it (mean plus hard
bound). A port should expose an $N_{\mathrm{step}}$ histogram counter from
day one; it is the single most informative health signal of the tier.

\subsection{Is tier 1 still necessary at long context, and does recall
depend on it?}\label{sec:tier1need}
Yes. The $99.8\%$ trigger rate is a miss \emph{frequency}, not a workload
share: at $S{=}512$k tier 1 supplies over $99\%$ of the attended rows
(${\approx}16{,}384$ kept vs ${\approx}120$ fetched per step), so tier 1 is
the hit path and recall the miss path of a cache hierarchy. Recall is also
defined in terms of tier 1 three times over (the fire criterion competes
against the kept-tier maximum, the entropy gate reads the kept-tier
softmax, and calibration defines hardness relative to the partition); the
no-competition counterfactual collapses to similarity top-$k$ (recall@4
$=0.489$ or $99.5\%$ fire rate). Full argument: the defense document
(\texttt{docs/defense/defense.tex}, Q1).

\subsection{Relation to ArkVale}
ArkVale (Chen et al., 2024) is the closest prior design --- evict,
summarize, recall on demand. This tier departs from it on every axis that
matters: individual \emph{rows} with an exact 64-dim summand plus a rank-$r$
sketch and a residual certificate, rather than lossy per-\emph{page}
bounding digests (a digest cannot be tight under RoPE --- measured
layer-uniformity $0.672$ NoPE vs.\ $0.024$ RoPE --- so ArkVale's regime is
structurally the loose one); admission by \emph{certified competition}
against the query's own kept-tier maximum with a self-calibrated guarantee,
rather than digest-similarity ranking; and query-\emph{independent}
eviction, so the partition exists before any query does. The two designs
also degenerate in dual directions under pressure (ArkVale: quality, as
pages fall off the budget; here: worst-case attended-set size, never
quality, since fetched rows are exact), and every recall-path optimization
above is NoPE-enabled with a RoPE counterfactual. The full five-point
comparison, the NoPE enablers, memory placement (no paging), and the
duality argument: the defense document
(\texttt{docs/defense/defense.tex}, Q2).


\subsection{The fetch cap: an optional deployment optimization (not part of the algorithm)}
\textbf{The serving default is uncapped} ($\mathrm{topj}=-1$): fetch the
whole fired set. A cap must be \emph{stated explicitly} by the operator; the
recommended value is $\mathrm{topj}=16$, which buys the bounded-fetch
guarantee
\[
\text{fetch} \;\le\; \mathrm{topj} \times H \ \text{rows / step / layer}
\]
(worst-case observed fire $3{,}983$ rows $\to$ capped $296$; a $13\times$
tightening) at measured-zero quality cost: generation bit-identical,
$\Delta$bpb $= 0.04$\,pp inside the pre-registered neutral band, needle
recovery identical. Rationale for explicitness: the cap is the one knob
trading recall completeness against a worst-case latency/PCIe bound, and the
engineer setting it should know which side they are buying.

%======================================================================
\section{Parameters}\label{sec:params}

\textbf{One quality knob.} The deployment exposes exactly one quality
parameter, the recall target $\tau=0.90$; the rest of the trigger
calibration is a closed-form chain, not tuned:
\[
\alpha=\tfrac{1-\tau}{2},\quad
\tau_s=\tfrac{\tau}{1-\alpha},\quad
z_p=z_{(k)}\;\text{with}\;k=\lceil(n{+}1)\tau_s\rceil,\quad
\mathrm{min\_hard}=\lceil\tfrac{\tau_s}{1-\tau_s}\rceil=18,
\]
plus an $8\!\to\!64$ adaptive calibration window and a $Z_{\max}=8$ clamp.

\textbf{Structural constants} (not tuned per deployment; each with a stated
reason):
\begin{center}\small
\begin{tabular}{@{}llp{7.1cm}@{}}
\toprule
constant & value & rationale / provenance\\
\midrule
$\rho$ & $1/32$ & compression ratio; the only ratio measured flat in context length (an experimental axis, not a runtime knob)\\
$\kappa$ & 16 & low-pass bandwidth of Eq.~\eqref{eq:sigma}; counts frequency bins, so a fixed block window keeps its meaning scale-invariant\\
$\mathrm{CLOSE\_BLOCK}$ & 4096 & the $\sigma$ transform window and close granularity; anchors $\kappa$'s scale invariance and makes $\sigma$ incremental\\
$r$ (sketch) & 64 & index rank; minimizes per-step scan bytes (the archive is GPU-resident, so fetches are free and larger $r$ only buys a lower fire rate)\\
sinks & 4 & always-kept head rows\\
unclosed tail & $<4096$ & the always-attended remainder since the last close; absorbs the recent-window role (a fixed $W{=}256$ window survives only as a debug checker's scale)\\
\bottomrule
\end{tabular}
\end{center}

\textbf{Optional optimization.} \texttt{topj} (fetch cap) defaults to $-1$
(uncapped, the deployed form); a bounded fetch (16 recommended) is an
optional worst-case-latency optimization, adopted from ArkVale's budgeted
admission and applied to the certified scores.

Operational values (partition, prefill chunking, backends, transport tuning)
are deployment-specific; the reference deployment records them, one file,
each entry carrying its value, rationale, and provenance, in \texttt{config/recommended.json} of the
experiment repository. \textbf{Fool-proof rule}: a default may do the whole
thing or nothing, never a fraction nobody typed --- which is why the cap
defaults off and algorithm knobs are typed explicitly at launch.

%======================================================================
\section{Decode-time block close and rebalance}\label{sec:close}

Tier-1 selection is not a one-shot prefill event. The decode loop closes one
block per $B=4096$ decoded tokens (\texttt{while seq\_len - closed >= B},
catching up multiple blocks if needed); each close, per (request, MLA layer):
\begin{enumerate}[itemsep=1pt]
\item compute $\sigma$ for the newly closed block (Eq.~\eqref{eq:sigma},
fused projection form) and append it to the immutable per-block record;
\item \textbf{global rebalance}: re-run Eq.~\eqref{eq:select} over all
$\mathrm{closed}$ rows with the grown quota $m=\max(1,\mathrm{round}(\rho\cdot\mathrm{closed}))$
--- rows move in either direction (a kept row whose $\sigma$ no longer makes
the cut is demoted to the archive);
\item rewrite the kept table as $[\text{selected closed rows} \;|\;
\text{unclosed tail}]$; the tail then grows one row per step until the next
close;
\item \textbf{refresh the recall index membership only}: extend the
closed-prefix operand caches (sketch/residual tables are append-only over
closed rows) and re-derive the archive/kept index selection. The calibrated
thresholds ($z$, $\theta_g$) are \emph{not} refit --- the conformal
certificate is sound for any archive membership, so it survives rebalance.
\end{enumerate}
The close runs synchronously off the captured graph (${\sim}1$\,ms per
$4096$ steps, amortized to noise), and bumps a pack epoch that makes the
in-graph buffers re-sync in place (contents, never addresses --- the decode
graph is not recaptured for a close).

%======================================================================
\section{Cost model and validation targets}\label{sec:cost}

Throughout, \emph{k} denotes $1024$ tokens (4k${=}4096$, 64k${=}65536$, 512k${=}524288$).

\textbf{Per-token attention volume is not a fixed budget} --- a property
classic sparse attention does not have, so it must enter the port's latency
contract explicitly. The tier-1 attended count is \emph{proportional}
($m + |\mathrm{tail}|$ with $m = \mathrm{round}(\rho T)$; a constant-$m$
schedule exists via \texttt{m\_fixed} but the proportional schedule is the
only one measured quality-flat in context length), and the recall tier makes
the per-step count \emph{query-dependent}: the fired set ranges from zero to
a measured worst case of $3{,}983$ rows. The algorithm itself is uncapped; the optional cap hard-bounds the
per-step fetch at $\mathrm{topj}\times H$ (the non-degeneration
guarantee) as a deployment optimization. Uncapped, the theoretical worst
case is full fire: that step attends the entire cache, i.e.\ naive MLA plus the scan
overhead (quality unharmed --- the fetched rows are exact; steps are
independent); the measured worst fire of $3{,}983$ rows is an observation,
not a bound. Kernel consequence: size the launch grid for
the maximum and rely on empty-split correctness
(Section~\ref{sec:kernel}); never re-derive the grid from the per-step
count in graph mode.

All validation targets in this section are \emph{tier-1 ablation}
figures. They are an ablation floor, not a deployable baseline: without the
recall tier, quality collapses at long context ($S \gtrsim 256$k), so
tier-1-only throughput must never be quoted as the configuration recall is
``slower than''. The only legitimate speed comparison is recall-on against
the uncompressed baseline.

\paragraph{Honest cost model, and where the speedup lives.} Decode step time
decomposes as
\[
t(S) \;=\; C_0 \;+\; T_{\mathrm{net}} \;+\; A(S),
\]
with $C_0$ the context-independent remainder (MoE experts, linear-attention
layers, sampling, host), $T_{\mathrm{net}}$ the pipeline-parallel transfer
cost (zero on a single device), and $A(S)$ the attention term --- the ONLY
term this method can shrink. Under the dense baseline
$A_{\mathrm{dense}}(S) = k\,S$ with $k$ the per-row read cost
($1152$\,B/row at bf16). Under VestigeKV,
\[
A_{\mathrm{vk}}(S) \;=\; \underbrace{k\,\rho S}_{\text{kept attention}}
\;+\; \underbrace{k_s (1-\rho) S}_{\text{archive scan}} \;+\; R,
\]
where $k_s$ is the scan's per-row cost ($260$\,B/row with the shipped
16-bit index: side bf16, sketch fp16, residual scalar fp32) and $R$ a
small fixed residual (the fused seven-kernel pipeline's per-step engineering
adder, ${\sim}68\,\mu$s by kernel-level profiling, batch-independent). Three consequences, stated
plainly:

\begin{enumerate}
\item \textbf{The advantage is speed, not memory.} Both tiers are
GPU-resident and the partition invariant deletes no row: the full latent
pool stays in device memory and the wrapper adds index buffers on top.
Memory savings would require host-offloading the archive (the $r{=}192$
variant of Section~\ref{sec:params}), which this implementation does not do.
\item \textbf{The $S$-slope does not vanish; it shrinks by
$\rho + k_s/k \approx 0.26$ with the shipped 16-bit index.} The scan must
touch every archived row every step --- that is what the recall guarantee
costs --- so the method cuts the attention slope to roughly a quarter
rather than eliminating it. (The fire-set stability gate that licensed the
16-bit scan operands has since passed; see Section~\ref{sec:params}.)
\item \textbf{Short-context speedup is bounded by attention's share of the
step.} On the reference deployment (two nodes over ethernet,
$T_{\mathrm{net}} \approx 2.0$\,ms measured, $C_0 \approx 5.5$\,ms,
$A_{\mathrm{dense}} \approx 0.57$\,ms at $S{=}64$k by measured-slope
calibration), attention is $7\%$ of the step at $64$k, so no selection
scheme --- this one or any other --- can win much there. Below ${\sim}32$k
the method is strictly at a disadvantage: the recall pipeline's fixed cost
remains while its saving vanishes, which is exactly the regime the optional
dense fallback (above) exists to skip in production. The reported curves run
with that fallback disabled, so the left edge of every figure shows this
cost undisguised.
\end{enumerate}

\textbf{This is a long-context method.} Measured on the reference
deployment (final fused build, bs=1 continuous decode), the vestigekv
step-time slope is ${\sim}0.5$\,ns/token against dense $6.3$\,ns/token;
the crossover sits at ${\sim}48$k and the end-to-end ratio grows with $S$:
$1.07\times$ (128k), $1.20\times$ (272k), $1.39\times$ (496k), extrapolating
toward the attention-only asymptote as $S$ grows (the saved term is linear
while $C_0 + T_{\mathrm{net}}$ is constant). The measured slope is
\emph{below} the byte account because the dense arm's own long-$S$ slope
exceeds its bandwidth floor (split-count growth and length-scaled kernel
inefficiencies), which vestigekv sidesteps by holding its attended set near
$\rho S$. The crossover is deployment-dependent in the honest direction:
the reference testbed's $C_0$ (MoE with 3B active parameters) and
$T_{\mathrm{net}}$ (ethernet pipeline) are LARGE, which suppresses the
ratio; a single-device or NVLink deployment moves the crossover earlier and
lifts every entry. Each arm's curve is collected within a single launch so
the fixed launch-state terms cancel; cross-arm claims read the same
statistic (median per 4k-token bucket) on both, and above ${\sim}240$k the
server-side decode log is the per-token authority.

Throughput (64k prefill + 4k decode, batch sweep): the attention share of
the step grows with batch (attention cannot batch while the shared weights
amortize), so the vestigekv advantage grows with it, against the dense
\texttt{triton} baseline: within launch-noise of even at bs=1 (attention is
a few percent of the step there), then $+7.7\%$ at bs=8 and $+8.3\%$ at
bs=12. The per-step engineering adder of the fused pipeline is
${\sim}68\,\mu$s (kernel-level profiling), batch-independent.

%======================================================================
\section{Validation protocol}
\label{sec:validate}

Four gate groups, each run before any number is trusted (full detail and
thresholds: the defense document, \texttt{docs/defense/defense.tex}, Q3):
\begin{enumerate}[itemsep=2pt]
\item \textbf{Plumbing parity}: kill-switch (wrapper disabled) and full-arm
(selection keeps \emph{all} rows) outputs must match stock call-for-call /
bit-for-bit; every fused kernel keeps its pre-fusion implementation as a
bit-parity oracle, and the fused pipeline's greedy output must be
token-identical to the unfused one.
\item \textbf{Selector equivalence}: $\sigma$ and $\keep$ must match the
reference implementation --- bit-exactly on the fallback path,
decision-exactly (identical keep masks) for the fused kernels
(\texttt{engine/test/manual/test\_vestigekv\_equiv.py}, self-contained).
\item \textbf{Quality gates} (uncapped, $\rho{=}1/32$): needle retrieval
1.00 at $8\times$--$32\times$ / 8k--65k; teacher-forced $\Delta$bpb vs the
uncompressed engine $\le 0.15\%$; recall at $1/128$ compression reaching
1.00 in the serving loop with the ${\sim}512$k necessity signature
($99.8\%$ trigger rate) reproducing.
\item \textbf{Regression-net discipline}: every test's guard must be
demonstrated to \emph{fail} on the broken form before it is trusted; the
16-bit scan index additionally passes a fire-set stability gate (Jaccard
$\ge 0.99$ vs the fp32 arm, each arm calibrated independently).
\end{enumerate}

%======================================================================
\section{Observed defect classes and mitigations}
\label{sec:pitfalls}

The port's postmortem surfaces twelve recurring defect classes; the four
structural ones are slot-keyed state going stale on slot reuse
(a re-used slot must not inherit the previous request's kept tables),
selection work landing inside the decode window (all $\sigma$ work belongs
at prefill / block close), per-request host loops on the step path
(vectorized away), and graph-replay reality (a backend's
\texttt{forward\_*} python does not run on replay --- per-step state must
live in the fixed-address buffers the metadata hooks touch). The rest are
measurement-hygiene classes: hidden device syncs in capture, padded-lane
overruns, arch-specific JIT of unrelated ops, configuration-matrix
isolation before attribution, host-side cost invisible to microbenchmarks,
benchmark warmup / launch-variance discipline, the second-graph launch tax,
and the capacity-grid launch floor. Each class, the bug it produced, and
the mitigation: the defense document
(\texttt{docs/defense/defense.tex}, Q4).

\vspace{1em}
\noindent\textbf{Provenance.} Selection math and quality numbers: the
pre-registered VestigeKV experimental record (PREREG series; RoPE control
measured on DeepSeek-V2-Lite). Engineering contract and performance:
the production port to standard \texttt{sglang}
(branch \texttt{vestigekv}), dual-GPU pipeline deployment of
Kimi Linear 48B-A3B; all changes outside the attention backend are
enumerated in that repository's port audit. This document froze against
commit \texttt{e44b914a7} (v2: selection-rule, gate-quantile, calibration
and close-semantics sections re-frozen against the code; see
\texttt{DIFFS\_zh.md} for the v1$\to$v2 corrections).

\end{document}
